\section{Mathematical details}
\subsection{Block reduction}
\label{sec:block_proof}

We use the component weights and attachment convention from
Section~\ref{sec:condensation}.

\begin{lemma}[Normalization of leaf weights]
\label{lem:normalization}
For each component, $\sum_{x\in S_a}\omega_a(x)=1$.
\end{lemma}
\begin{proof}
A singleton has weight one.
At a binary root, every leaf weight in each child component is multiplied by $1/2$.
Induction over the component therefore gives total mass
$\tfrac12\cdot1+\tfrac12\cdot1=1$.
\end{proof}

\paragraph{Proof of Theorem~\ref{thm:clade_condense}.}

\begin{proof}
Separate the unordered leaf pairs into pairs within a component and pairs between components.
For $x,y\in S_a$, the path stays inside $S_a$ and is unchanged by rearranging the skeleton.
The sum of all such terms is
\[
 C=\sum_a\sum_{\substack{x<y\\x,y\in S_a}}
      D_{xy}\,2^{-q_{T_G}(x,y)},
\]
which is independent of $G$.

For $x\in S_a$ and $y\in S_b$ with $a\ne b$, the path contains the branching vertices inside $S_a$, those on the skeleton path, and those inside $S_b$.
These sets of vertices are disjoint, giving
\[
 q_{T_G}(x,y)=h_a(x)+q_G(a,b)+h_b(y).
\]
Therefore its weight factors as
$2^{-q_{T_G}(x,y)}=\omega_a(x)\,2^{-q_G(a,b)}\,\omega_b(y)$.
Summing between-component terms gives
\begin{align*}
 \sum_{a<b}\sum_{\substack{x\in S_a\\y\in S_b}}
 D_{xy}\,2^{-q_{T_G}(x,y)}
 &=\sum_{a<b}2^{-q_G(a,b)}
      \sum_{\substack{x\in S_a\\y\in S_b}}\omega_a(x)\omega_b(y)D_{xy}\\
 &=L(G;\widetilde D).
\end{align*}
Adding the within-component constant proves Equation~\eqref{eq:bme_condense_decomp}.
Subtracting two such decompositions cancels $C$ and proves Equation~\eqref{eq:condensed_delta}.
\end{proof}

\subsection{Ordered representation}
\label{sec:gauge}

For an ordered vector, initialize $r[x]=x$ and process $i=N-2,\ldots,0$.
Join components $r[v_i]$ and $r[i+1]$ at node $2N-2-i$, then replace $r[v_i]$ by this node.
The active components have distinct minimum labels $0,\ldots,i+1$ before each merge, giving an $O(N)$ decoder.

Relabeling uses Penn et al.'s linear-time construction~\cite[Lemmas~0.9--0.11]{penn2023gradme}.
Starting with root label zero and counter $t=1$, visit internal nodes parent before child.
Each passes its label to one child, labels the other $t$, records the inherited label as $v_{t-1}$, and increments $t$.
Our first orientation lets the child with smaller minimum leaf index inherit; the second reverses that choice.
Both preserve the tree and its BME score when both axes of $D$ are permuted by the inverse leaf-label map.
The construction guarantees representational coverage after relabeling; optimizer reachability depends on its search policy.
For unrooted tree-to-vector conversion, insert the artificial root on the edge incident to leaf zero.

\subsection{Coordinate replacements are SPR moves}
\label{sec:coord_spr}

For an ordered $v$ and a leaf $\ell\geq i+1$, follow the chain $\ell_0=\ell$, $\ell_{s+1}=v_{\ell_s-1}$ while $\ell_s>i+1$.
Each step uses a coordinate with index greater than $i$, and the chain strictly decreases.
Let $P_i$ be the set of leaves whose chain stops at $i+1$; it contains $i+1$ and does not depend on $v_i$.

\begin{lemma}[Coordinate neighbors are SPR neighbors]
\label{lem:coord_spr}
Let $v\in\mathcal V_N^{\mathrm{cf}}$, $0\leq i\leq N-2$, and $0\leq j\leq2i$.
Then $T(v^{(i\leftarrow j)})$ either equals $T(v)$ or is obtained from $T(v)$ by one SPR move that prunes the subtree on the leaves $P_i$.
\end{lemma}
\begin{proof}
Read $w=v^{(i\leftarrow j)}$ by leaf addition~\cite{penn2024phylo2vec}.
Starting from leaf $0$, coordinate $w_k$ inserts leaf $k+1$ on one of the $2k+1$ edges of the current rooted tree on leaves $0,\ldots,k$, including the edge above its root, and a value $w_k\leq k$ selects the pendant edge of leaf $w_k$.
Only step $i$ depends on $j$.
It inserts leaf $i+1$ on some edge, creating a vertex $p$.
Every later step $k>i$ has $w_k=v_k\leq k$ and inserts leaf $k+1$ on the pendant edge of leaf $v_k$.
The new leaf therefore lies on the same side of $p$ as leaf $v_k$.
By induction on $k$, the leaves on the side of $p$ that contains $i+1$ are exactly $P_i$.
The insertions among these leaves build a subtree whose topology does not depend on $j$.
Removing $P_i$ and suppressing $p$ leaves the tree $T^-$ produced by steps $0,\ldots,i-1$ and the insertions of leaves outside $P_i$, which also does not depend on $j$.
Thus every $T(w)$ is $T^-$ with the fixed subtree on $P_i$ attached at one edge of the unrooted $T^-$.
If leaf $i+1$ is inserted above the root, that edge is the one through the root of $T^-$, which becomes a single edge when the artificial root is suppressed.
Taking $j=v_i$ recovers $T(v)$.
Two such trees differ only in the attachment edge of $P_i$, so they are equal or related by one SPR move.
\end{proof}

The lemma bounds the coordinate neighborhood by the SPR neighborhood.
It does not show the reverse inclusion.
The mapping below identifies the attachment edge for each $j$.

\subsection{Mapping coordinate values to regraft edges}
\label{sec:mapping}

The ordered decoder numbers leaves $0,\ldots,N-1$ and internal nodes $N,\ldots,2N-2$.
The last index is the artificial root.
Suppressing it leaves a computational root $R=2N-3$ and the top edge joining $R$ to the artificial root's other child, denoted $c_{\mathrm{other}}$.
Index ordinary edges by their child endpoints.

For row $i$, let $E_i[u]$ be the score change from pruning the component associated with leaf $i+1$ and regrafting it on the edge indexed by $u$.
The half-tree SPR traversal supplies these values.
Let $b_i(j)$ be the root of the component containing leaf $j$ after the merges with coordinate index greater than $i$.
Then the edge-to-coordinate map is
\begin{equation}
 \Delta_{i,j}=
 \begin{cases}
 E_i[b_i(j)], &0\leq j\leq i,\\[2pt]
 E_i[(2N-2)-2i+j], &i+1\leq j\leq2i-2,\\[2pt]
 E_i[c_{\mathrm{other}}], &j\in\{2i-1,2i\}\cap[i+1,2i].
 \end{cases}
 \label{eq:full_mapping}
\end{equation}
The first case identifies the component on the corresponding pendant edge of the prefix tree.
In the ordered decoder's node numbering, the remaining non-root prefix edges occupy consecutive indices $2N-1-i,\ldots,2N-4$, giving the affine map in the second case.
The last two rooted edge choices collapse onto the same unrooted top edge.
At $i=1$, the middle case is empty and $j=1$ belongs to the first case.
Only $j=2$ takes the collapsed-root value.
Coordinate zero has no nontrivial alternative.

On a CPU, the array $b_i$ can be maintained incrementally as the merge order advances.
On a GPU, independent rows reconstruct their mappings using the latest eligible merge for each minimum leaf.
When several merges have the same minimum leaf, only the latest eligible merge may write that destination.
This avoids concurrent writes that would otherwise depend on thread scheduling.

\subsection{Parallel evaluation and reduced storage}
\label{sec:parallel_details}

After initialization and the required transpose, Equation~\eqref{eq:bottom_up} has independent columns.
For a non-root ancestor $u$ of $w$, let $\operatorname{sib}_u(w)$ be the child of $u$ whose subtree excludes $w$.
The complementary direction is computed from the root down:
\begin{equation}
 A(u,w)=\tfrac12\{A(\operatorname{sib}_u(w),w)
                         +A(\operatorname{parent}(u),w)\}.
 \label{eq:top_down}
\end{equation}
The degree-three root is initialized using its two branches away from $w$.

For reduced storage, retain only internal-to-leaf distances.
These initialize the leaf entries of an internal row $w$; a leaf row instead starts from $D$.
Bottom-up averaging fills internal entries, and Equation~\eqref{eq:top_down} overwrites ancestors of $w$ after the root initialization.
Set the diagonal to zero.
For a vertex $u$ in a disjoint component, the bottom-up calculation follows directly by linearity.
For a descendant $u$ of $w$, the initial leaf entries already use the component at $w$ facing away from $u$, so the same averaging is valid.
Ancestors require the complementary direction supplied by the overwrite.
Thus each row costs $O(N+\operatorname{depth}(w))=O(N)$ to reconstruct.
A constant number of reconstructions per pruned component retains $O(N^2)$ sweep work while reducing the balanced-distance table from approximately $4N^2$ to $N^2$ entries.

\section{Search settings}
\label{sec:search_policies}

Algorithms~\ref{alg:block_search} and~\ref{alg:coordinate_search} use the following settings, with scores evaluated in the recorded component labeling.
The CPU solver uses sweep budget $I=500$, escape budget $E=2$, and $\eta=10^{-12}(|\ell|+1)$ at score $\ell$.

\paragraph{Block construction.}
For requested size $3\leq K_\star<N$, expose $w$ consecutive leaves in the current depth-first order, with
$w=\min\{K_\star-2,\max\{4,\lfloor2K_\star/3\rfloor\}\}$.
Block $b$ starts at $b\max\{1,\lfloor w/2\rfloor\}\bmod N$.
Condense maximal inactive subtrees, then split the largest such component until the skeleton reaches $K_\star$ leaves or no split remains possible.
The resulting size $K$ can differ from the target $K_\star$.
Rebuild the leaf order, components, and outside information from the current tree for every block.
Hierarchical schedules combine backbone and local windows with staggered or larger windows crossing component boundaries.
Full-tree search uses $K=N$ singleton components.

\paragraph{Row selection and joint acceptance.}
$\proc{RankMin}$ selects $j_i\in\arg\min_{0\leq j\leq2i}\Delta_{i,j}$, retains changes below $-\eta$, and sorts them increasingly, resolving ties by deterministic traversal and row order.
Let $d_{\min}<0$ be the best single change.
$\proc{BatchStep}$ greedily packs at most eight rows, excluding shared marked pruned-subtree nodes and attachment endpoints.
After the first row, require $\Delta_{i,j_i}<0.05d_{\min}$, and test a batch only if the sum is below $1.05d_{\min}$.
Each trial starts from the same pre-trial encoding $v$ and is accepted only if its complete score satisfies
$L(T(v');M_\pi)<\ell+d_{\min}-\eta$.
On rejection, retry successively halved ranked prefixes down to two rows; if none passes, use the best single replacement.
The predicted sum only screens candidates.

\paragraph{Completion and escapes.}
$\proc{Decode}$ restores component identities and suppresses the artificial root.
$\proc{ExtraSPR}$ and $\proc{RemainingSPR}$ return an encoding in the current labeling or $\bot$ if no change is below $-\eta$.
After an unsuccessful completion scan, an escape is permitted only if a subsequent sweep remains.
For each row, retain the smallest change in $(\eta,0.05)$ and sort these increasingly.
Attempt $e\in\{1,2\}$ combines the first and $(e+1)$th candidates if both exist and their sum is below $0.02$; otherwise it returns $\bot$.
The thresholds are in objective units.
Escapes can increase the working score, so the solver returns the best scored tree, including the final trial.

\paragraph{GPU policy.}
The GPU solver packs SPR moves with disjoint marked paths, rescores each combined trial on the next sweep, and retains the best scored skeleton over a bounded number of sweeps.
It tries reversed labeling when no improving row remains, without CPU batch-halving or full SPR completion.
The best trial is checked against the entry skeleton before committing a block.
Intermediate trial scores can increase.
By default (\texttt{--init sample}), sequence refinement uses sample-and-place initialization ($n_0=\min(100{,}000,\lfloor N/10\rfloor)$) followed by up to $R_{\max}=8$ outer V-cycle rounds with early stopping when a round improves the BME objective by less than $10^{-6}$ (\texttt{P2V\_STOP\_ZERO=1}, \texttt{P2V\_STOP\_TOL=1e-6}).
For sequence refinement above $2{,}048$ leaves, inner skeleton proposals cap component leaf-weight lists at $64$ entries (\texttt{max\_prof\_leaves = 64}); in default scaling runs, window commits use the capped surrogate directly (\texttt{--eta 0}), while an optional lazy exact verification mode (\texttt{--eta} $>0$) constructs full leaf profiles when $\Delta L_{\mathrm{capped}}>-\eta_{\mathrm{exact}}$ and checks Equation~\eqref{eq:changed_pairs} in FP64.
Following topology refinement, Corollary~\ref{cor:branch_lengths} estimates all $2N-3$ FastME/BME branch lengths on GPU using top-$32$ inside and outside leaf-weight profiles per internal node.

\clearpage
\section{Additional experimental results}
\label{sec:additional_results}

\subsection{Sequence-refinement scaling tables}
\label{sec:jc_tables}

Tables~\ref{tab:rnasim_scaling}, \ref{tab:jc10k_scaling}, and~\ref{tab:jc1k_scaling} provide the complete numerical measurements corresponding to Figures~\ref{fig:rnasim_scaling}, \ref{fig:jc_scaling_4panel}, and~\ref{fig:seqlen_effect} on simulated RNASim ($L=1{,}596$) and Jukes--Cantor ($L=10{,}000$ and $1{,}000$) alignments.
At $N\in\{1{,}000{,}000,10{,}000{,}000\}$, where the ground-truth trees are constructed by expanding each leaf of the $100{,}000$-taxon reference tree into $10$- or $100$-leaf clades, we report both the overall normalized RF across all $N-3$ splits and (in brackets) the normalized RF restricted to the $99{,}997$ non-trivial reference backbone splits.
We also report the \texttt{Phylo2Vec-Minibatch (RandStart)} ablation, which starts directly from the same uniform random tree passed to \texttt{-intree} CPU baselines without sample-and-place initialization.

\begin{table*}[!htbp]
\centering
\caption{Sequence-to-tree scaling on simulated \textbf{RNASim} alignments ($L=1{,}596$ sites) from $N=1{,}000$ to $10{,}000{,}000$ leaves, comparing \textbf{Phylo2Vec-Minibatch} (sample-and-place initialization + multi-round V-cycle refinement + GPU branch lengths) and \textbf{Phylo2Vec-Minibatch (RandStart)} (uniform random start) against \texttt{DIPPER 0.1.5} (\textit{de novo}) and CPU baselines initialized from the same uniform random trees ($100.0\%$ initial RF). Entries for $N\leq100{,}000$ report $\text{mean}\pm\text{SD}$ over $R=5$ independent replicates; at $N\ge 10^6$, brackets denote RF (\%) on the $99{,}997$ reference backbone splits.}
\label{tab:rnasim_scaling}
\scriptsize
\setlength{\tabcolsep}{3.5pt}
\resizebox{\linewidth}{!}{%
\begin{tabular}{@{}rrlrrrrrr@{}}
\toprule
$N$ & $R$ & Method & RF (\%) & Time (s) & $\Delta\text{BME}$ & Host RAM (MB) & VRAM (MB) & Hardware\\
\midrule
1,000 & 5 & Uniform Random Init & $99.97 \pm 0.04$ & $0.00 \pm 0.00$ & $0.00 \pm 0.00$ & --- & --- & ---\\
1,000 & 5 & DIPPER 0.1.5 (GPU, \textit{de novo}) & $20.98 \pm 0.82$ & $0.47 \pm 0.07$ & --- & $121.6 \pm 0.5$ & $392.0$ & 64 CPU + $1\times$ A6000\\
1,000 & 5 & FastME 2.1.6 (\texttt{-s -u}) & $18.05 \pm 0.84$ & $114.08 \pm 9.45$ & --- & $123.3 \pm 1.2$ & --- & 64 CPU\\
1,000 & 5 & FastTreeMP 2.2 (\texttt{-noml -intree}) & $19.70 \pm 3.12$ & $11.94 \pm 0.20$ & --- & $133.5 \pm 2.9$ & --- & 64 CPU\\
1,000 & 5 & VeryFastTree 4.0 (\texttt{-noml -intree}) & $18.62 \pm 0.89$ & $8.04 \pm 0.24$ & --- & $131.6 \pm 2.2$ & --- & 64 CPU\\
1,000 & 5 & FastTreeMP 2.2 (ME+ML \texttt{-intree}) & $\mathbf{14.42 \pm 0.47}$ & $18.30 \pm 0.58$ & --- & $133.5 \pm 2.9$ & --- & 64 CPU\\
1,000 & 5 & VeryFastTree 4.0 (ME+ML \texttt{-intree}) & $14.62 \pm 0.65$ & $9.42 \pm 0.27$ & --- & $133.5 \pm 2.9$ & --- & 64 CPU\\
1,000 & 5 & Phylo2Vec-Minibatch (RandStart) & $18.31 \pm 1.19$ & $0.05 \pm 0.00$ & $-144.10 \pm 1.45$ & $376.2 \pm 0.0$ & $688.8$ & 64 CPU + $8\times$ H100\\
1,000 & 5 & \textbf{Phylo2Vec-Minibatch} & $18.29 \pm 0.71$ & $\mathbf{0.03 \pm 0.00}$ & $-38.69 \pm 1.48$ & $376.2 \pm 0.0$ & $688.8$ & 64 CPU + $8\times$ H100\\
\midrule
5,000 & 5 & Uniform Random Init & $100.00 \pm 0.00$ & $0.00 \pm 0.00$ & $0.00 \pm 0.00$ & --- & --- & ---\\
5,000 & 5 & DIPPER 0.1.5 (GPU, \textit{de novo}) & $18.82 \pm 0.77$ & $2.07 \pm 0.06$ & --- & $133.0 \pm 0.0$ & $580.0$ & 64 CPU + $1\times$ A6000\\
5,000 & 5 & FastTreeMP 2.2 (\texttt{-noml -intree}) & $45.55 \pm 9.45$ & $70.38 \pm 1.91$ & --- & $290.4 \pm 0.5$ & --- & 64 CPU\\
5,000 & 5 & VeryFastTree 4.0 (\texttt{-noml -intree}) & $43.17 \pm 9.67$ & $42.96 \pm 0.86$ & --- & $317.6 \pm 7.4$ & --- & 64 CPU\\
5,000 & 5 & FastTreeMP 2.2 (ME+ML \texttt{-intree}) & $\mathbf{13.62 \pm 1.33}$ & $117.48 \pm 5.45$ & --- & $292.2 \pm 1.8$ & --- & 64 CPU\\
5,000 & 5 & VeryFastTree 4.0 (ME+ML \texttt{-intree}) & $14.85 \pm 1.41$ & $50.22 \pm 1.55$ & --- & $318.6 \pm 8.0$ & --- & 64 CPU\\
5,000 & 5 & Phylo2Vec-Minibatch (RandStart) & $14.80 \pm 0.62$ & $0.54 \pm 0.05$ & $-869.75 \pm 6.89$ & $379.4 \pm 0.7$ & $694.8$ & 64 CPU + $8\times$ H100\\
5,000 & 5 & \textbf{Phylo2Vec-Minibatch} & $14.90 \pm 0.59$ & $\mathbf{0.35 \pm 0.02}$ & $-175.19 \pm 1.46$ & $379.4 \pm 0.7$ & $694.8$ & 64 CPU + $8\times$ H100\\
\midrule
10,000 & 5 & Uniform Random Init & $100.00 \pm 0.00$ & $0.01 \pm 0.01$ & $0.00 \pm 0.00$ & --- & --- & ---\\
10,000 & 5 & DIPPER 0.1.5 (GPU, \textit{de novo}) & $17.78 \pm 0.38$ & $8.90 \pm 0.10$ & --- & $145.8 \pm 0.8$ & $1156.0$ & 64 CPU + $1\times$ A6000\\
10,000 & 5 & FastTreeMP 2.2 (\texttt{-noml -intree}) & $64.34 \pm 7.47$ & $155.42 \pm 5.57$ & --- & $570.1 \pm 1.3$ & --- & 64 CPU\\
10,000 & 5 & VeryFastTree 4.0 (\texttt{-noml -intree}) & $58.97 \pm 6.78$ & $90.56 \pm 3.09$ & --- & $542.3 \pm 7.1$ & --- & 64 CPU\\
10,000 & 5 & FastTreeMP 2.2 (ME+ML \texttt{-intree}) & $14.47 \pm 1.33$ & $280.93 \pm 11.76$ & --- & $571.8 \pm 1.3$ & --- & 64 CPU\\
10,000 & 5 & VeryFastTree 4.0 (ME+ML \texttt{-intree}) & $14.29 \pm 0.59$ & $103.90 \pm 3.22$ & --- & $545.0 \pm 6.8$ & --- & 64 CPU\\
10,000 & 5 & Phylo2Vec-Minibatch (RandStart) & $\mathbf{13.65 \pm 0.40}$ & $1.79 \pm 0.91$ & $-1868.34 \pm 9.98$ & $384.6 \pm 0.7$ & $698.8$ & 64 CPU + $8\times$ H100\\
10,000 & 5 & \textbf{Phylo2Vec-Minibatch} & $13.67 \pm 0.37$ & $\mathbf{0.95 \pm 0.10}$ & $-342.55 \pm 6.80$ & $384.8 \pm 0.5$ & $698.8$ & 64 CPU + $8\times$ H100\\
\midrule
50,000 & 5 & Uniform Random Init & $100.00 \pm 0.00$ & $0.04 \pm 0.00$ & $0.00 \pm 0.00$ & --- & --- & ---\\
50,000 & 5 & DIPPER 0.1.5 (GPU, \textit{de novo}) & $19.44 \pm 0.17$ & $56.84 \pm 0.72$ & --- & $246.6 \pm 0.8$ & $446.0$ & 64 CPU + $1\times$ A6000\\
50,000 & 5 & FastTreeMP 2.2 (\texttt{-noml -intree}) & $89.30 \pm 1.13$ & $874.43 \pm 24.09$ & --- & $2835.9 \pm 0.6$ & --- & 64 CPU\\
50,000 & 5 & VeryFastTree 4.0 (\texttt{-noml -intree}) & $88.09 \pm 1.29$ & $486.95 \pm 38.28$ & --- & $2405.4 \pm 15.5$ & --- & 64 CPU\\
50,000 & 5 & Phylo2Vec-Minibatch (RandStart) & $\mathbf{9.91 \pm 0.15}$ & $5.97 \pm 0.60$ & $-10943.39 \pm 12.09$ & $428.5 \pm 2.3$ & $744.8$ & 64 CPU + $8\times$ H100\\
50,000 & 5 & \textbf{Phylo2Vec-Minibatch} & $9.95 \pm 0.12$ & $\mathbf{5.30 \pm 0.41}$ & $-1698.05 \pm 14.85$ & $430.9 \pm 1.9$ & $744.8$ & 64 CPU + $8\times$ H100\\
\midrule
100,000 & 5 & Uniform Random Init & $100.00 \pm 0.00$ & $0.10 \pm 0.00$ & $0.00 \pm 0.00$ & --- & --- & ---\\
100,000 & 5 & DIPPER 0.1.5 (GPU, \textit{de novo}) & $16.80 \pm 0.19$ & $216.26 \pm 0.22$ & --- & $386.5 \pm 2.8$ & $512.0$ & 64 CPU + $1\times$ A6000\\
100,000 & 5 & FastTreeMP 2.2 (\texttt{-noml -intree}) & $92.62 \pm 0.59$ & $1853.11 \pm 67.03$ & --- & $5677.2 \pm 2.3$ & --- & 64 CPU\\
100,000 & 5 & VeryFastTree 4.0 (\texttt{-noml -intree}) & $92.73 \pm 1.37$ & $1024.52 \pm 37.34$ & --- & $4651.7 \pm 24.7$ & --- & 64 CPU\\
100,000 & 5 & Phylo2Vec-Minibatch (RandStart) & $8.75 \pm 0.05$ & $12.53 \pm 0.64$ & $-23262.69 \pm 20.40$ & $470.4 \pm 0.8$ & $792.8$ & 64 CPU + $8\times$ H100\\
100,000 & 5 & \textbf{Phylo2Vec-Minibatch} & $\mathbf{8.74 \pm 0.06}$ & $\mathbf{8.34 \pm 0.49}$ & $-3326.69 \pm 23.90$ & $470.4 \pm 0.7$ & $792.8$ & 64 CPU + $8\times$ H100\\
\midrule
1,000,000 & 1 & Uniform Random Init & $100.00$ & $0.00$ & $0.00$ & --- & --- & ---\\
1,000,000 & 1 & DIPPER 0.1.5 (GPU, \textit{de novo}) & $11.25$ & $978.23$ & --- & $4189.6$ & $1871.8$ & 64 CPU + $1\times$ A6000\\
1,000,000 & 1 & Phylo2Vec-Minibatch (RandStart) & $1.01$ [$9.09$] & $108.06$ & $-306187.06$ & $\mathbf{1050.8}$ & $\mathbf{1656.8}$ & 64 CPU + $8\times$ H100\\
1,000,000 & 1 & \textbf{Phylo2Vec-Minibatch} & $\mathbf{1.00}$ [$\mathbf{9.08}$] & $\mathbf{36.61}$ & $-24068.02$ & $1531.0$ & $\mathbf{1656.8}$ & 64 CPU + $8\times$ H100\\
\midrule
10,000,000 & 1 & Uniform Random Init & $100.00$ & $0.00$ & $0.00$ & --- & --- & ---\\
10,000,000 & 1 & \textbf{Phylo2Vec-Minibatch} & $\mathbf{0.12}$ [$\mathbf{9.15}$] & $\mathbf{430.48}$ & $-458725.25$ & $\mathbf{7092.8}$ & $10240.8$ & 64 CPU + $8\times$ H100\\
\bottomrule
\end{tabular}%
}
\end{table*}

\begin{table*}[!htbp]
\centering
\caption{Sequence-to-tree scaling on simulated Jukes--Cantor alignments with $L=10{,}000$ sites (\texttt{JC10k}) from $N=1{,}000$ to $10{,}000{,}000$ leaves, comparing \textbf{Phylo2Vec-Minibatch} (sample-and-place initialization + multi-round V-cycle refinement + GPU branch lengths) and \textbf{Phylo2Vec-Minibatch (RandStart)} (uniform random start) against \texttt{DIPPER 0.1.5} (\textit{de novo}) and CPU baselines initialized from the same uniform random trees ($100.0\%$ initial RF). Entries for $N\leq100{,}000$ report $\text{mean}\pm\text{SD}$ over $R=5$ independent replicates; at $N\ge 10^6$, brackets denote RF (\%) on the $99{,}997$ reference backbone splits.}
\label{tab:jc10k_scaling}
\scriptsize
\setlength{\tabcolsep}{3.5pt}
\resizebox{\linewidth}{!}{%
\begin{tabular}{@{}rrlrrrrrr@{}}
\toprule
$N$ & $R$ & Method & RF (\%) & Time (s) & $\Delta\text{BME}$ & Host RAM (MB) & VRAM (MB) & Hardware\\
\midrule
1,000 & 5 & Uniform Random Init & $99.97 \pm 0.04$ & $0.00 \pm 0.00$ & $0.00 \pm 0.00$ & --- & --- & ---\\
1,000 & 5 & DIPPER 0.1.5 (GPU, \textit{de novo}) & $6.60 \pm 0.48$ & $0.59 \pm 0.10$ & --- & $134.2 \pm 1.0$ & $398.0$ & 64 CPU + $1\times$ A6000\\
1,000 & 5 & FastME 2.1.6 (\texttt{-s -u}) & $4.25 \pm 0.57$ & $110.76 \pm 4.85$ & --- & $138.0 \pm 2.5$ & --- & 64 CPU\\
1,000 & 5 & FastTreeMP 2.2 (\texttt{-noml -intree}) & $4.25 \pm 0.57$ & $68.20 \pm 1.47$ & --- & $345.0 \pm 0.6$ & --- & 64 CPU\\
1,000 & 5 & VeryFastTree 4.0 (\texttt{-noml -intree}) & $4.99 \pm 1.36$ & $47.15 \pm 1.77$ & --- & $599.4 \pm 37.4$ & --- & 64 CPU\\
1,000 & 5 & FastTreeMP 2.2 (ME+ML \texttt{-intree}) & $\mathbf{2.61 \pm 0.52}$ & $91.23 \pm 1.30$ & --- & $395.4 \pm 3.9$ & --- & 64 CPU\\
1,000 & 5 & VeryFastTree 4.0 (ME+ML \texttt{-intree}) & $2.63 \pm 0.67$ & $52.11 \pm 1.45$ & --- & $603.1 \pm 44.1$ & --- & 64 CPU\\
1,000 & 5 & Phylo2Vec-Minibatch (RandStart) & $4.29 \pm 0.58$ & $0.05 \pm 0.01$ & $-132.79 \pm 1.63$ & $381.5 \pm 0.0$ & $696.8$ & 64 CPU + $8\times$ H100\\
1,000 & 5 & \textbf{Phylo2Vec-Minibatch} & $4.25 \pm 0.57$ & $\mathbf{0.03 \pm 0.00}$ & $-21.88 \pm 1.78$ & $381.2 \pm 0.7$ & $696.8$ & 64 CPU + $8\times$ H100\\
\midrule
5,000 & 5 & Uniform Random Init & $100.00 \pm 0.00$ & $0.00 \pm 0.00$ & $0.00 \pm 0.00$ & --- & --- & ---\\
5,000 & 5 & DIPPER 0.1.5 (GPU, \textit{de novo}) & $6.01 \pm 0.31$ & $2.32 \pm 0.10$ & --- & $193.8 \pm 0.1$ & $600.0$ & 64 CPU + $1\times$ A6000\\
5,000 & 5 & FastTreeMP 2.2 (\texttt{-noml -intree}) & $22.02 \pm 8.00$ & $380.93 \pm 5.04$ & --- & $1676.2 \pm 3.6$ & --- & 64 CPU\\
5,000 & 5 & VeryFastTree 4.0 (\texttt{-noml -intree}) & $9.27 \pm 2.30$ & $238.48 \pm 5.14$ & --- & $1666.2 \pm 47.3$ & --- & 64 CPU\\
5,000 & 5 & FastTreeMP 2.2 (ME+ML \texttt{-intree}) & $2.04 \pm 0.31$ & $541.68 \pm 22.96$ & --- & $1742.7 \pm 4.4$ & --- & 64 CPU\\
5,000 & 5 & VeryFastTree 4.0 (ME+ML \texttt{-intree}) & $\mathbf{1.98 \pm 0.36}$ & $256.49 \pm 5.46$ & --- & $1649.8 \pm 45.5$ & --- & 64 CPU\\
5,000 & 5 & Phylo2Vec-Minibatch (RandStart) & $3.33 \pm 0.12$ & $1.08 \pm 0.03$ & $-738.28 \pm 9.79$ & $405.3 \pm 1.4$ & $720.8$ & 64 CPU + $8\times$ H100\\
5,000 & 5 & \textbf{Phylo2Vec-Minibatch} & $3.36 \pm 0.13$ & $\mathbf{0.83 \pm 0.01}$ & $-82.17 \pm 1.03$ & $404.6 \pm 1.7$ & $720.8$ & 64 CPU + $8\times$ H100\\
\midrule
10,000 & 5 & Uniform Random Init & $100.00 \pm 0.00$ & $0.01 \pm 0.01$ & $0.00 \pm 0.00$ & --- & --- & ---\\
10,000 & 5 & DIPPER 0.1.5 (GPU, \textit{de novo}) & $6.21 \pm 0.11$ & $9.31 \pm 0.05$ & --- & $266.6 \pm 0.8$ & $1196.0$ & 64 CPU + $1\times$ A6000\\
10,000 & 5 & FastTreeMP 2.2 (\texttt{-noml -intree}) & $33.51 \pm 13.07$ & $823.06 \pm 11.99$ & --- & $3345.3 \pm 9.4$ & --- & 64 CPU\\
10,000 & 5 & VeryFastTree 4.0 (\texttt{-noml -intree}) & $33.55 \pm 11.07$ & $475.03 \pm 13.05$ & --- & $2766.5 \pm 14.1$ & --- & 64 CPU\\
10,000 & 5 & FastTreeMP 2.2 (ME+ML \texttt{-intree}) & $2.00 \pm 0.14$ & $1230.72 \pm 91.55$ & --- & $3411.9 \pm 9.2$ & --- & 64 CPU\\
10,000 & 5 & VeryFastTree 4.0 (ME+ML \texttt{-intree}) & $\mathbf{1.92 \pm 0.06}$ & $525.70 \pm 20.07$ & --- & $2770.2 \pm 26.4$ & --- & 64 CPU\\
10,000 & 5 & Phylo2Vec-Minibatch (RandStart) & $3.07 \pm 0.08$ & $3.33 \pm 0.06$ & $-1542.67 \pm 15.40$ & $435.1 \pm 1.4$ & $750.8$ & 64 CPU + $8\times$ H100\\
10,000 & 5 & \textbf{Phylo2Vec-Minibatch} & $3.05 \pm 0.10$ & $\mathbf{2.73 \pm 0.47}$ & $-149.89 \pm 2.55$ & $435.8 \pm 0.0$ & $750.8$ & 64 CPU + $8\times$ H100\\
\midrule
50,000 & 5 & Uniform Random Init & $100.00 \pm 0.00$ & $0.04 \pm 0.00$ & $0.00 \pm 0.00$ & --- & --- & ---\\
50,000 & 5 & DIPPER 0.1.5 (GPU, \textit{de novo}) & $4.94 \pm 0.18$ & $65.60 \pm 1.75$ & --- & $1031.8 \pm 8.2$ & $646.0$ & 64 CPU + $1\times$ A6000\\
50,000 & 4 & FastTreeMP 2.2 (\texttt{-noml -intree}) & $74.25 \pm 4.51$ & $4525.52 \pm 115.67$ & --- & $16673.0 \pm 19.7$ & --- & 64 CPU\\
50,000 & 5 & VeryFastTree 4.0 (\texttt{-noml -intree}) & $70.03 \pm 2.04$ & $2506.54 \pm 103.01$ & --- & $12062.2 \pm 61.7$ & --- & 64 CPU\\
50,000 & 5 & Phylo2Vec-Minibatch (RandStart) & $\mathbf{2.56 \pm 0.11}$ & $21.31 \pm 3.80$ & $-8408.38 \pm 56.20$ & $630.7 \pm 1.4$ & $1002.8$ & 64 CPU + $8\times$ H100\\
50,000 & 5 & \textbf{Phylo2Vec-Minibatch} & $2.57 \pm 0.11$ & $\mathbf{14.17 \pm 0.05}$ & $-640.75 \pm 5.15$ & $\mathbf{629.6 \pm 0.7}$ & $1002.8$ & 64 CPU + $8\times$ H100\\
\midrule
100,000 & 5 & Uniform Random Init & $100.00 \pm 0.00$ & $0.10 \pm 0.00$ & $0.00 \pm 0.00$ & --- & --- & ---\\
100,000 & 5 & DIPPER 0.1.5 (GPU, \textit{de novo}) & $4.38 \pm 0.14$ & $249.05 \pm 1.33$ & --- & $1995.0 \pm 3.9$ & $912.0$ & 64 CPU + $1\times$ A6000\\
100,000 & 5 & VeryFastTree 4.0 (\texttt{-noml -intree}) & $84.88 \pm 2.64$ & $5588.08 \pm 300.03$ & --- & $23767.2 \pm 110.4$ & --- & 64 CPU\\
100,000 & 5 & Phylo2Vec-Minibatch (RandStart) & $2.43 \pm 0.12$ & $37.28 \pm 3.98$ & $-17306.97 \pm 122.28$ & $631.1 \pm 0.9$ & $1308.8$ & 64 CPU + $8\times$ H100\\
100,000 & 5 & \textbf{Phylo2Vec-Minibatch} & $\mathbf{2.42 \pm 0.12}$ & $\mathbf{26.08 \pm 0.31}$ & $-1186.91 \pm 10.02$ & $\mathbf{630.0 \pm 1.1}$ & $1308.8$ & 64 CPU + $8\times$ H100\\
\midrule
1,000,000 & 1 & Uniform Random Init & $100.00$ & $0.00$ & $0.00$ & --- & --- & ---\\
1,000,000 & 1 & DIPPER 0.1.5 (GPU, \textit{de novo}) & $5.64$ & $2592.75$ & --- & $24225.3$ & $17747.6$ & 64 CPU + $1\times$ A6000\\
1,000,000 & 1 & Phylo2Vec-Minibatch (RandStart) & $1.18$ [$8.31$] & $338.78$ & $-201273.28$ & $\mathbf{1046.0}$ & $\mathbf{6806.8}$ & 64 CPU + $8\times$ H100\\
1,000,000 & 1 & \textbf{Phylo2Vec-Minibatch} & $\mathbf{0.32}$ [$\mathbf{2.80}$] & $\mathbf{145.94}$ & $-11863.02$ & $1534.2$ & $\mathbf{6806.8}$ & 64 CPU + $8\times$ H100\\
\midrule
10,000,000 & 1 & Uniform Random Init & $100.00$ & $0.00$ & $0.00$ & --- & --- & ---\\
10,000,000 & 1 & \textbf{Phylo2Vec-Minibatch} & $\mathbf{0.03}$ [$\mathbf{2.87}$] & $\mathbf{1161.52}$ & $-312968.56$ & $\mathbf{7077.0}$ & $61738.8$ & 64 CPU + $8\times$ H100\\
\bottomrule
\end{tabular}%
}
\end{table*}

\begin{table*}[!htbp]
\centering
\caption{Sequence-length sensitivity at $N=10{,}000$ leaves: simulated Jukes--Cantor alignments with $L=1{,}000$ sites (\texttt{JC1k}) vs.\ $L=10{,}000$ sites (\texttt{JC10k}) over $R=5$ independent replicates ($\text{mean}\pm\text{SD}$).}
\label{tab:jc1k_scaling}
\scriptsize
\setlength{\tabcolsep}{3.5pt}
\begin{tabular*}{\linewidth}{@{\extracolsep{\fill}}lrrrrrr@{}}
\toprule
Method & \multicolumn{3}{c}{$L=1{,}000$ sites (\texttt{JC1k}, $R=5$)} & \multicolumn{3}{c}{$L=10{,}000$ sites (\texttt{JC10k}, $R=5$)}\\
\cmidrule(lr){2-4}\cmidrule(lr){5-7}
 & Norm.\ RF (\%) & Time (s) & $\Delta\text{BME}$ & Norm.\ RF (\%) & Time (s) & $\Delta\text{BME}$\\
\midrule
Uniform Random Init & $100.00 \pm 0.00$ & $0.01 \pm 0.01$ & $0.00 \pm 0.00$ & $100.00 \pm 0.00$ & $0.01 \pm 0.01$ & $0.00 \pm 0.00$\\
DIPPER 0.1.5 (GPU, \textit{de novo}) & $22.06 \pm 0.26$ & $8.74 \pm 0.08$ & --- & $6.21 \pm 0.11$ & $9.31 \pm 0.05$ & ---\\
FastTreeMP 2.2 (\texttt{-noml -intree}) & $47.52 \pm 10.80$ & $85.00 \pm 3.06$ & --- & $33.51 \pm 13.07$ & $823.06 \pm 11.99$ & ---\\
VeryFastTree 4.0 (\texttt{-noml -intree}) & $41.36 \pm 5.67$ & $50.23 \pm 2.15$ & --- & $33.55 \pm 11.07$ & $475.03 \pm 13.05$ & ---\\
FastTreeMP 2.2 (ME+ML \texttt{-intree}) & $12.18 \pm 0.80$ & $129.78 \pm 5.78$ & --- & $2.00 \pm 0.14$ & $1230.72 \pm 91.55$ & ---\\
VeryFastTree 4.0 (ME+ML \texttt{-intree}) & $\mathbf{11.87 \pm 0.63}$ & $55.72 \pm 2.27$ & --- & $\mathbf{1.92 \pm 0.06}$ & $525.70 \pm 20.07$ & ---\\
Phylo2Vec-Minibatch (RandStart) & $13.25 \pm 0.23$ & $1.00 \pm 0.07$ & $-1561.84 \pm 26.90$ & $3.07 \pm 0.08$ & $3.33 \pm 0.06$ & $-1542.67 \pm 15.40$\\
\textbf{Phylo2Vec-Minibatch} & $13.27 \pm 0.33$ & $\mathbf{0.59 \pm 0.01}$ & $-159.13 \pm 4.69$ & $3.05 \pm 0.10$ & $\mathbf{2.73 \pm 0.47}$ & $-149.89 \pm 2.55$\\
\bottomrule
\end{tabular*}
\end{table*}

\begin{table*}[!htbp]
\centering
\caption{$O(N)$ GPU FastME/BME branch-length estimation scaling (Corollary~\ref{cor:branch_lengths}, top-$32$ dyadic profiles per internal node) from $N=1{,}000$ to $10{,}000{,}000$ leaves ($19{,}999{,}997$ unrooted edges) on $8\times\text{NVIDIA H100}$ GPUs (64 CPU threads).}
\label{tab:bme_bl_scaling}
\scriptsize
\setlength{\tabcolsep}{4.5pt}
\begin{tabular*}{\linewidth}{@{\extracolsep{\fill}}rrrrrrrr@{}}
\toprule
$N$ & Unrooted edges ($2N-3$) & \multicolumn{3}{c}{\textbf{RNASim} ($L=1{,}596$ sites, $W=25$)} & \multicolumn{3}{c}{\textbf{JC10k} ($L=10{,}000$ sites, $W=160$)}\\
\cmidrule(lr){3-5}\cmidrule(lr){6-8}
 &  & CPU profiles (s) & GPU kernel (s) & Total (s) & CPU profiles (s) & GPU kernel (s) & Total (s)\\
\midrule
1,000 & 1,997 & 0.0005 & 0.0025 & \textbf{0.0029} & 0.0004 & 0.0039 & \textbf{0.0043}\\
5,000 & 9,997 & 0.0025 & 0.0029 & \textbf{0.0054} & 0.0014 & 0.0047 & \textbf{0.0061}\\
10,000 & 19,997 & 0.0023 & 0.0048 & \textbf{0.0071} & 0.0044 & 0.0077 & \textbf{0.0121}\\
50,000 & 99,997 & 0.0325 & 0.0076 & \textbf{0.0400} & 0.0335 & 0.0129 & \textbf{0.0464}\\
100,000 & 199,997 & 0.0732 & 0.0273 & \textbf{0.1004} & 0.0652 & 0.0219 & \textbf{0.0872}\\
1,000,000 & 1,999,997 & 0.4779 & 0.1210 & \textbf{0.5989} & 0.4649 & 0.1972 & \textbf{0.6621}\\
10,000,000 & 19,999,997 & 4.4016 & 1.3923 & \textbf{5.7939} & 5.8504 & 2.0030 & \textbf{7.8534}\\
\bottomrule
\end{tabular*}
\end{table*}

\begin{table*}[!htbp]
\centering
\caption{Pandemic-scale SARS-CoV-2 phylogenetic inference against \texttt{CMAPLE}~\cite{lytrong2024cmaple} and \texttt{MAPLE}~\cite{demaio2023maple} across the official simulated ($N\in\{100\text{K},200\text{K},500\text{K}\}$, $L=30{,}000$~bp) and empirical GISAID ($N=10\text{K}\to1\text{M}$, $L=29{,}891$~bp) benchmark suites. For simulated datasets, we report all-splits normalized RF error (\texttt{RF\_all}) and resolved-branch RF error on branches with expected length $\ge0.5/L$ (\texttt{RF\_0.5}) and $\ge1.0/L$ (\texttt{RF\_1.0}). \texttt{Phylo2Vec-Minibatch} times include sample-and-place initialization ($n_0=N/10$), V-cycle refinement, and $O(N)$ GPU branch lengths.}
\label{tab:cmaple_comparison}
\scriptsize
\setlength{\tabcolsep}{3.0pt}
\resizebox{\textwidth}{!}{%
\begin{tabular}{@{}llrrrrrrrrrrrrrr@{}}
\toprule
Dataset & $N$ & \multicolumn{4}{c}{\textbf{CMAPLE} (ML)} & \multicolumn{4}{c}{\textbf{MAPLE} (ML)} & \multicolumn{4}{c}{\textbf{Phylo2Vec-Minibatch} ($8\times$ H100)} & \multicolumn{2}{c}{Speedup vs.}\\
\cmidrule(lr){3-6}\cmidrule(lr){7-10}\cmidrule(lr){11-14}\cmidrule(lr){15-16}
 & & Time (s) & RAM (GB) & \texttt{RF\_all} & \texttt{RF\_1.0} & Time (s) & RAM (GB) & \texttt{RF\_all} & \texttt{RF\_1.0} & Time (s) & RAM (GB) & \texttt{RF\_all} & \texttt{RF\_1.0} & CMAPLE & MAPLE\\
\midrule
Sim ($\bar\ell=5\cdot10^{-5}$) & 100K & 1,010.7 & 2.87 & 24.29\% & 10.08\% & 4,169.0 & 5.63 & 24.32\% & 10.12\% & \textbf{74.8} & \textbf{0.99} & 24.71\% & 10.35\% & $13.5\times$ & $55.7\times$\\
Sim ($\bar\ell=5\cdot10^{-5}$) & 200K & 3,866.4 & 5.80 & 24.40\% & 10.03\% & 10,886.0 & 11.46 & 24.42\% & 10.03\% & \textbf{109.7} & \textbf{0.99} & 24.64\% & 10.21\% & $35.2\times$ & $99.2\times$\\
Sim ($\bar\ell=5\cdot10^{-5}$) & 500K & 8,559.2 & 14.00 & 23.96\% & 9.81\% & 25,653.0 & 26.92 & 23.95\% & 9.80\% & \textbf{214.5} & \textbf{1.14} & 24.32\% & 9.94\% & $39.9\times$ & $119.6\times$\\
\midrule
Sim ($\bar\ell=10^{-4}$) & 100K & 1,189.3 & 6.94 & 11.90\% & 5.86\% & 5,543.0 & 14.82 & 11.91\% & 5.86\% & \textbf{76.0} & \textbf{0.99} & 12.00\% & 5.90\% & $15.6\times$ & $72.9\times$\\
Sim ($\bar\ell=10^{-4}$) & 200K & 2,470.7 & 14.44 & 11.67\% & 5.74\% & 7,954.0 & 30.38 & 11.67\% & 5.74\% & \textbf{112.1} & \textbf{0.99} & 11.75\% & 5.80\% & $22.1\times$ & $71.0\times$\\
Sim ($\bar\ell=10^{-4}$) & 500K & 5,511.9 & 37.41 & 11.68\% & 5.71\% & 23,771.0 & 79.47 & 11.68\% & 5.71\% & \textbf{212.9} & \textbf{1.15} & 11.75\% & 5.77\% & $25.9\times$ & $111.6\times$\\
\midrule
Sim ($\bar\ell=1.5\cdot10^{-4}$) & 100K & 1,267.3 & 12.58 & 7.26\% & 4.23\% & 8,781.0 & 28.21 & 7.27\% & 4.24\% & \textbf{76.3} & \textbf{0.99} & 7.32\% & 4.27\% & $16.6\times$ & $115.0\times$\\
Sim ($\bar\ell=1.5\cdot10^{-4}$) & 200K & 3,443.7 & 24.20 & 7.17\% & 4.16\% & 23,829.0 & 55.99 & 7.17\% & 4.17\% & \textbf{114.3} & \textbf{0.99} & 7.26\% & 4.20\% & $30.1\times$ & $208.5\times$\\
Sim ($\bar\ell=1.5\cdot10^{-4}$) & 500K & 8,550.3 & 64.19 & 7.12\% & 4.08\% & 51,608.0 & 145.89 & 7.12\% & 4.08\% & \textbf{213.2} & \textbf{1.15} & 7.16\% & 4.10\% & $40.1\times$ & $242.1\times$\\
\midrule
GISAID Empirical & 10K & 496.8 & --- & --- & --- & --- & --- & --- & --- & \textbf{9.9} & \textbf{0.91} & --- & --- & $50.1\times$ & ---\\
GISAID Empirical & 50K & 7,524.5 & --- & --- & --- & --- & --- & --- & --- & \textbf{69.7} & \textbf{0.99} & --- & --- & $108.0\times$ & ---\\
GISAID Empirical & 100K & 36,602.3 & --- & --- & --- & --- & --- & --- & --- & \textbf{125.6} & \textbf{0.99} & --- & --- & $291.4\times$ & ---\\
GISAID Empirical & 200K & 70,076.9 & --- & --- & --- & --- & --- & --- & --- & \textbf{170.3} & \textbf{0.99} & --- & --- & $411.6\times$ & ---\\
GISAID Empirical & 500K & 215,055.0 & --- & --- & --- & --- & --- & --- & --- & \textbf{333.1} & \textbf{1.10} & --- & --- & $645.5\times$ & ---\\
GISAID Empirical & 1M & 943,582.2 & --- & --- & --- & $60{,}424.7^\dagger$ & --- & --- & --- & \textbf{609.0} & \textbf{1.40} & --- & --- & $\mathbf{1{,}549.5\times}$ & $99.2\times^\dagger$\\
\bottomrule
\multicolumn{16}{@{}l}{\footnotesize $^\dagger$At $N=1\text{M}$ empirical GISAID, \texttt{MAPLE} column reports placement-only \texttt{CMAPLE -no-tree-search} ($60{,}424.7$~s $=16.78$~h).}
\end{tabular}%
}
\end{table*}

\clearpage
\subsection{Numerical verification and small-tree CPU optimization}
\label{sec:cpu_tables}

Table~\ref{tab:validation} compares coordinate, condensed, and branch-length calculations with independent direct BME references.
Tables~\ref{tab:bench_noiseless} and~\ref{tab:bench_noisy} report single-core CPU optimization comparisons against the in-memory FastME BME-SPR reference (FME) on Apple M5 across additive and noisy distance matrices ($N=20$ to $200$).

\begin{table}[!htbp]
\centering
\caption{Agreement with independent BME calculations. Errors are maximum absolute discrepancies over the stated tests.}
\label{tab:validation}
\small
\begin{tabular}{@{}p{0.39\linewidth}p{0.32\linewidth}r@{}}
\toprule
Calculation & Test scope & Maximum error\\
\midrule
CPU coordinate differences & 36,000 changes, $N=20$ & $3.553\times10^{-14}$\\
Ordered-tree enumeration & 28,488 changes, $N=4$--7 & $2.776\times10^{-15}$\\
Relabeling, first orientation & 200 trees, $N=25$ & $1.776\times10^{-14}$\\
Relabeling, reversed orientation & Same 200 trees & $2.665\times10^{-14}$\\
Condensed SPR differences & 16,231 changes, $N=40$, $K=16$ & $1.564\times10^{-13}$\\
Pauplin tree-length identity $\sum_e\hat\ell(e)=L(T;D)$ & $N=20$--1,000 ($1{,}997$ edges) & $2.27\times10^{-13}$\\
FP32 CUDA gradient & Reported sizes $N=100$--4,000 & $1.11\times10^{-6}$\\
FP32 CUDA BME branch lengths (uncapped) & $1{,}997$ edges, $N=1{,}000$ & $3.74\times10^{-6}$\\
FP32 CUDA BME branch lengths (top-$64$ cap, MAE) & $1{,}997$ edges, $N=1{,}000$ & $6.13\times10^{-5}$\\
\bottomrule
\end{tabular}
\end{table}

\begin{table}[!htbp]
\centering\small
\setlength{\tabcolsep}{4pt}
\caption{Additive distance comparison on Apple M5.}
\label{tab:bench_noiseless}
\par\smallskip
\textit{Objective and topology}\par\smallskip
\begin{tabular}{@{}rrrrrr@{}}
\toprule
$N$ & FME $L$ & P2V $L$ & Win/tie/loss & FME RF & P2V RF\\
\midrule
$20$ & $11.12739$ & $11.12739$ & \textbf{0 / 20 / 0} & $0.0000$ & $0.0000$\\
$50$ & $29.42572$ & $29.42572$ & \textbf{0 / 20 / 0} & $0.0000$ & $0.0000$\\
$100$ & $59.52693$ & $59.52693$ & \textbf{0 / 10 / 0} & $0.0000$ & $0.0000$\\
$150$ & $89.04365$ & $89.04365$ & \textbf{0 / 5 / 0} & $0.0000$ & $0.0000$\\
$200$ & $116.32398$ & $116.32398$ & \textbf{0 / 5 / 0} & $0.0000$ & $0.0000$\\
\bottomrule
\end{tabular}
\par\smallskip
\textit{Computational cost}\par\smallskip
\begin{tabular}{@{}rrrrrrr@{}}
\toprule
$N$ & FME iters. & P2V evals. & FME ms & P2V ms & P2V med. [IQR] & Ratio\\
\midrule
$20$ & $13.4$ & $11.2$ & $0.12 \pm 0.02$ & $0.08 \pm 0.01$ & $0.07\,[0.02]$ & \textbf{1.54$\times$}\\
$50$ & $42.2$ & $19.4$ & $1.58 \pm 0.35$ & $0.57 \pm 0.13$ & $0.59\,[0.16]$ & \textbf{2.78$\times$}\\
$100$ & $97.7$ & $35.4$ & $9.05 \pm 0.84$ & $2.63 \pm 0.26$ & $2.65\,[0.44]$ & \textbf{3.44$\times$}\\
$150$ & $156.2$ & $49.4$ & $31.47 \pm 0.77$ & $7.92 \pm 0.86$ & $7.75\,[0.64]$ & \textbf{3.97$\times$}\\
$200$ & $211.0$ & $63.6$ & $75.46 \pm 2.42$ & $18.14 \pm 0.74$ & $18.05\,[1.17]$ & \textbf{4.16$\times$}\\
\bottomrule
\end{tabular}

\end{table}

\begin{table}[!htbp]
\centering\small
\setlength{\tabcolsep}{4pt}
\caption{Noisy distance comparison on Apple M5.}
\label{tab:bench_noisy}
\par\smallskip
\textit{Objective and topology}\par\smallskip
\begin{tabular}{@{}rrrrrr@{}}
\toprule
$N$ & FME $L$ & P2V $L$ & Win/tie/loss & FME RF & P2V RF\\
\midrule
$20$ & $11.19777$ & $11.19777$ & \textbf{0 / 20 / 0} & $0.0500$ & $0.0500$\\
$50$ & $29.25507$ & $29.25607$ & \textbf{1 / 18 / 1} & $0.0660$ & $0.0649$\\
$100$ & $59.25471$ & $59.25471$ & \textbf{0 / 10 / 0} & $0.0526$ & $0.0526$\\
$150$ & $92.14846$ & $92.14846$ & \textbf{0 / 5 / 0} & $0.0395$ & $0.0395$\\
$200$ & $120.94220$ & $120.94220$ & \textbf{1 / 4 / 0} & $0.0467$ & $0.0487$\\
\bottomrule
\end{tabular}
\par\smallskip
\textit{Computational cost}\par\smallskip
\begin{tabular}{@{}rrrrrrr@{}}
\toprule
$N$ & FME iters. & P2V evals. & FME ms & P2V ms & P2V med. [IQR] & Ratio\\
\midrule
$20$ & $14.4$ & $12.4$ & $0.07 \pm 0.01$ & $0.05 \pm 0.01$ & $0.04\,[0.02]$ & \textbf{1.60$\times$}\\
$50$ & $46.1$ & $25.9$ & $1.07 \pm 0.07$ & $0.49 \pm 0.09$ & $0.49\,[0.10]$ & \textbf{2.18$\times$}\\
$100$ & $102.6$ & $41.0$ & $10.62 \pm 2.53$ & $3.32 \pm 0.56$ & $3.18\,[0.52]$ & \textbf{3.19$\times$}\\
$150$ & $166.0$ & $61.0$ & $33.97 \pm 1.01$ & $10.15 \pm 1.35$ & $10.17\,[1.35]$ & \textbf{3.35$\times$}\\
$200$ & $235.8$ & $84.4$ & $85.16 \pm 2.06$ & $24.67 \pm 2.79$ & $24.51\,[4.48]$ & \textbf{3.45$\times$}\\
\bottomrule
\end{tabular}

\end{table}


\clearpage
\subsection{GPU evaluation}
\label{sec:scaling_tables}

Tables~\ref{tab:cuda_a6000_verify}--\ref{tab:h100_scaling} report coordinate-sweep and condensation costs at fixed trees.
Single-GPU speedups use different host CPU references.
In Table~\ref{tab:h100_scaling}, $2N^2$ is a rectangular-workload proxy; the number of valid coordinate entries is $(N-1)^2$.
Memory units follow the benchmark records.

For Table~\ref{tab:minibatch_scaling}, each skeleton contains 15,000 expanded leaves and 5,000 condensed components.
Epoch estimates multiply mean condensation-plus-sweep time by $\lceil N/15{,}000\rceil$ and exclude additional search work and optimization to convergence.
Validation discrepancies are observed errors on the benchmark checks.

\begin{table}[!htbp]
\centering\small
\setlength{\tabcolsep}{4pt}
\caption{Single-GPU gradient timings and agreement with the CPU reference. Times are milliseconds.}
\label{tab:cuda_a6000_verify}
\par\smallskip
\textit{RTX A6000 timings and reported maximum errors}\par\smallskip
\begin{tabular}{@{}rrrrr@{}}
\toprule
$N$ & A6000 full & Ratio & A6000 tiled & Max. error\\
\midrule
$100$ & $0.061$ & \textbf{8.44$\times$} & $0.078$ & $3.68\text{e-}07$\\
$250$ & $0.186$ & \textbf{15.92$\times$} & $0.207$ & $6.22\text{e-}07$\\
$500$ & $0.382$ & \textbf{32.95$\times$} & $0.511$ & $6.48\text{e-}07$\\
$1{,}000$ & $1.035$ & \textbf{33.69$\times$} & $1.988$ & $9.29\text{e-}07$\\
$2{,}000$ & $3.851$ & \textbf{40.15$\times$} & $7.288$ & $1.11\text{e-}06$\\
$4{,}000$ & $11.053$ & \textbf{62.48$\times$} & $43.124$ & $1.07\text{e-}06$\\
\bottomrule
\end{tabular}
\par\smallskip
\textit{H100 timings}\par\smallskip
\begin{tabular}{@{}rrrrrrr@{}}
\toprule
$N$ & CPU ref. & H100 avg. & H100 SPR & H100 total & Ratio & H100 tiled\\
\midrule
$100$ & $0.249$ & $0.044$ & $0.027$ & $0.071$ & \textbf{3.51$\times$} & $0.095$\\
$250$ & $1.964$ & $0.116$ & $0.048$ & $0.164$ & \textbf{11.98$\times$} & $0.205$\\
$500$ & $8.176$ & $0.227$ & $0.077$ & $0.304$ & \textbf{26.89$\times$} & $0.460$\\
$1{,}000$ & $32.336$ & $0.661$ & $0.105$ & $0.766$ & \textbf{42.21$\times$} & $1.030$\\
$2{,}000$ & $154.091$ & $1.557$ & $0.726$ & $2.283$ & \textbf{67.50$\times$} & $4.273$\\
$4{,}000$ & $701.375$ & $3.565$ & $2.654$ & $6.219$ & \textbf{112.78$\times$} & $22.327$\\
\bottomrule
\end{tabular}
\par\medskip\begin{minipage}{\linewidth}\footnotesize The maximum-error column is the reported combined check across the GPU paths. All listed checks pass their benchmark tolerance. CPU baselines differ between hosts.\end{minipage}
\end{table}

\begin{table}[!htbp]
\centering\small
\setlength{\tabcolsep}{4pt}
\caption{Embedding-based condensed sweeps with $K=20{,}000$ skeleton leaves. The epoch columns are batch-count times mean condensation-plus-gradient time.}
\label{tab:minibatch_scaling}
\par\smallskip
\textit{Two A6000 GPUs}\par\smallskip
\begin{tabular}{@{}rrrr@{}}
\toprule
$N$ & Batches & 2 A6000 cond.+grad. (ms) & Epoch (s)\\
\midrule
$1{,}000{,}000$ & $67$ & $4.45 + 284.44$ & $19.36$\\
$5{,}000{,}000$ & $334$ & $8.42 + 283.02$ & $97.34$\\
$10{,}000{,}000$ & $667$ & $17.29 + 281.58$ & $199.35$\\
\bottomrule
\end{tabular}
\par\smallskip
\textit{Eight H100 GPUs}\par\smallskip
\begin{tabular}{@{}rrrrrr@{}}
\toprule
$N$ & Cond. ms & Grad. ms & Epoch s & MB/GPU & Check\\
\midrule
$1{,}000{,}000$ & $2.69$ & $48.53$ & $\mathbf{3.43}$ & $7{,}793$ & \textbf{PASS} ($1.4\text{e-}05$)\\
$5{,}000{,}000$ & $4.86$ & $48.48$ & $\mathbf{17.82}$ & $7{,}824$ & \textbf{PASS} ($2.1\text{e-}06$)\\
$10{,}000{,}000$ & $10.20$ & $48.57$ & $\mathbf{39.20}$ & $7{,}862$ & \textbf{PASS} ($2.2\text{e-}06$)\\
\bottomrule
\end{tabular}
\par\medskip\begin{minipage}{\linewidth}\footnotesize This workload uses a bilinear embedding distance and is separate from sequence-distance refinement. The validation figures are observed discrepancies, not universal error bounds.\end{minipage}
\end{table}

\begin{table}[!htbp]
\centering\small
\setlength{\tabcolsep}{4pt}
\caption{Full-table and tiled multi-GPU sweeps. Times are milliseconds and memory is reported in MB.}
\label{tab:cuda_a6000_scaling}
\par\smallskip
\textit{A6000 timings and memory columns}\par\smallskip
\begin{tabular}{@{}rrrrr@{}}
\toprule
$N$ & 2 A6000 full & 2 A6000 tiled & Full memory & Tiled memory\\
\midrule
$1{,}000$ & $0.93$ & $5.75$ & $27$ & $39$\\
$2{,}000$ & $4.89$ & $13.90$ & $92$ & $93$\\
$4{,}000$ & $14.69$ & $56.85$ & $336$ & $247$\\
$10{,}000$ & $90.45$ & $249.59$ & $1{,}985$ & $1{,}076$\\
$20{,}000$ & $422.73$ & $1{,}212.27$ & $7{,}785$ & $3{,}677$\\
$40{,}000$ & \emph{OOM} & $5{,}628.62$ & $31{,}142$ & $13{,}614$\\
\bottomrule
\end{tabular}
\par\smallskip
\textit{Full-table H100 sweep}\par\smallskip
\begin{tabular}{@{}rrrr@{}}
\toprule
$N$ & 1 H100 [avg.+SPR] & 8 H100 & Ratio\\
\midrule
$1{,}000$ & $0.76\; [0.65 + 0.11]$ & $0.79$ & $0.96\times$\\
$2{,}000$ & $2.21\; [1.47 + 0.74]$ & $1.71$ & \textbf{1.29$\times$}\\
$4{,}000$ & $6.01\; [3.57 + 2.44]$ & $3.98$ & \textbf{1.51$\times$}\\
$10{,}000$ & $29.61\; [10.10 + 19.50]$ & $13.03$ & \textbf{2.27$\times$}\\
$20{,}000$ & $155.93\; [23.23 + 132.70]$ & $41.06$ & \textbf{3.80$\times$}\\
$40{,}000$ & $816.34\; [60.08 + 756.26]$ & $\mathbf{158.36}$ & \textbf{5.15$\times$}\\
\bottomrule
\end{tabular}
\par\smallskip
\textit{Tiled H100 sweep}\par\smallskip
\begin{tabular}{@{}rrrr@{}}
\toprule
$N$ & 1 H100 tiled & 8 H100 tiled & Top-1 check\\
\midrule
$1{,}000$ & $1.46$ & $4.47$ & \textbf{PASS} ($0.00\text{e+}00$)\\
$2{,}000$ & $6.14$ & $5.21$ & \textbf{PASS} ($2.38\text{e-}07$)\\
$4{,}000$ & $30.78$ & $10.98$ & \textbf{PASS} ($0.00\text{e+}00$)\\
$10{,}000$ & $250.69$ & $57.53$ & \textbf{PASS} ($0.00\text{e+}00$)\\
$20{,}000$ & $1{,}123.43$ & $244.13$ & \textbf{PASS} ($4.77\text{e-}07$)\\
$40{,}000$ & $4{,}768.31$ & $\mathbf{1{,}032.83}$ & \textbf{PASS} ($0.00\text{e+}00$)\\
\bottomrule
\end{tabular}

\end{table}

\begin{table}[!htbp]
\centering\small
\setlength{\tabcolsep}{4pt}
\caption{Distributed sweeps on eight H100 GPUs.}
\label{tab:h100_scaling}
\par\smallskip
\textit{Timing}\par\smallskip
\begin{tabular}{@{}rrrrr@{}}
\toprule
$N$ & Work proxy $2N^2$ & Distance ms & SPR ms & Total s\\
\midrule
$40{,}000$ & $3.20 \times 10^9$ & $514.54$ & $517.92$ & $\mathbf{1.032}$\\
$80{,}000$ & $1.28 \times 10^{10}$ & $2{,}112.82$ & $2{,}321.28$ & $\mathbf{4.434}$\\
$120{,}000$ & $2.88 \times 10^{10}$ & $4{,}785.75$ & $4{,}803.60$ & $\mathbf{9.589}$\\
$180{,}000$ & $6.48 \times 10^{10}$ & $10{,}848.67$ & $13{,}977.73$ & $\mathbf{24.826}$\\
$250{,}000$ & $1.25 \times 10^{11}$ & $39{,}981.55$ & $43{,}721.48$ & $\mathbf{83.703}$\\
$350{,}000$ & $2.45 \times 10^{11}$ & $76{,}240.05$ & $96{,}515.83$ & $\mathbf{172.756}$\\
$380{,}000$ & $2.89 \times 10^{11}$ & $89{,}992.00$ & $115{,}457.57$ & $\mathbf{205.450}$\\
\bottomrule
\end{tabular}
\par\smallskip
\textit{Memory and execution mode}\par\smallskip
\begin{tabular}{@{}rrrr@{}}
\toprule
$N$ & MB/GPU & Aggregate GB & Mode and check\\
\midrule
$40{,}000$ & $2{,}386$ & $18.6$ & \texttt{Full-Shard} [\textbf{PASS}]\\
$80{,}000$ & $7{,}824$ & $61.1$ & \texttt{Shard+Calib} [\textbf{PASS}]\\
$120{,}000$ & $16{,}314$ & $127.5$ & \texttt{Shard+Calib} [\textbf{PASS}]\\
$180{,}000$ & $34{,}770$ & $271.6$ & \texttt{Shard+Calib} [\textbf{PASS}]\\
$250{,}000$ & $62{,}296$ & $486.7$ & \texttt{Shard+Calib} [\textbf{PASS}]\\
$350{,}000$ & $62{,}180$ & $485.8$ & \texttt{StreamD+Calib} [\textbf{PASS}]\\
$380{,}000$ & $72{,}946$ & $569.9$ & \texttt{StreamD+Calib} [\textbf{PASS}]\\
\bottomrule
\end{tabular}
\par\medskip\begin{minipage}{\linewidth}\footnotesize At $N=400{,}000$, the reported requirement is 80,600 MB/GPU (629.7 GB aggregate), exceeding available device memory. Full-Shard and Shard+Calib retain distance shards. StreamD+Calib generates distance tiles as required. The work proxy is not a count of distinct moves.\end{minipage}
\end{table}

